\scanpage{54}
\mspage{10}

des graphes pondérés à repères, l'opération \deleted{s'explicite} se reflète par une opération correspondante au niveau de $\vec R$, par le ``gommage'' d'une partie $\vec R'$ stable par ${\unsure{\mathbb D|D}}_\infty$ i.e. par $\sigma_0,\sigma_1$

2) Gommage de cercles de découpage

On se donne maintenant une \deleted{ens. de composantes connexes de} $1$-sous-variété \inserted{(ouv. et) fermée} $\Sigma'$ de $\Sigma\setminus\partial X$, correspondant à un ensemble \deleted{\unsure{des}} de \unsure{comp.} conn. de $\unsure{\Sigma'|\Sigma}$ i.e. à une partie \inserted{$A'_0$} de $\pi_0(\Sigma')=A_0$ (arêtes ordinaires du graphe) --- et on remplace $\Sigma$ par $\Sigma\setminus\Sigma'$. Au niveau du graphe pondéré, cela se traduit par une opération \uline{de contraction} \uline{en un point} des arêtes correspondantes, d'où formation d'un nouveau graphe, et d'un \uline{morphisme} \uline{de générisation} (à l'opposé de ``spécialisation'') de l'ancien vers le nouveau. Il convient ici d'expliciter comment la pondération suit\ldots\ Au niveau des \inserted{sommets-}repères, il est entendu que les sommets repères se trouvant sur $\Sigma'$ sont gommés en $=$ temps --- donc, sur le graphe pondéré à repères, on fait sauter l'ens des repères qui sont au dessus des arêtes $\in A_{\unsure{0|o}}'$.
% archive sheet number: 53

